\section{Discussion and Limitations}

Three diagnostics were applied to a system whose aggregate score, MAcc
\num{0.866}, sits comfortably within the published range for this corpus, and
each isolates a different weakness that the headline number hides.

\emph{Unsupervised structure (k-means on the feature spaces).} In two of three
feature spaces, $k$-means recovers the acquisition site far better than the
diagnosis, so the class boundary is a thin supervised direction inside a geometry
organised by how and where the recording was made.

\emph{Per-site performance (best model, SVM--MFCC).} MAcc ranges from
\num{0.986} to \num{0.477} across sub-databases, with the pooled figure carried
by the single largest and easiest one; performance does not transfer across
acquisition conditions.

\emph{Attribution (Grad-CAM on the CNN).} The CNN's attention is genuinely
weight-dependent and mildly systole-preferring, but no more so on correct
decisions than on incorrect ones, so the tendency cannot be read as use of
systolic evidence.

None of the three establishes a mechanism on its own, so we ran the experiment
that ties them together (Section~\ref{sec:siteshortcut}): predicting the
sub-database directly from the features, and re-evaluating the diagnosis across
sites rather than pooled. It confirms the inference, with one qualification.
The site is recoverable at \num{0.647} balanced accuracy against
\num{0.167} chance, so the shortcut substrate is real; and pooled diagnosis MAcc
collapses from \num{0.840} to \num{0.457} under leave-one-site-out evaluation,
so almost none of the headline performance transfers across acquisition
conditions. The qualification is that a within-site signal does not vanish:
held at a fixed site, diagnosis MAcc is \num{0.70}--\num{0.76}, and even after
discounting two tiny or extreme-prevalence sites that inflate that average, the
large sites retain a modest \num{0.55}--\num{0.64}. The features are therefore
not \emph{only} a site fingerprint. The defensible summary is the measured one:
the aggregate metric is substantially, not entirely, a site-prevalence artefact,
and it overstates the transferable diagnostic signal by roughly the
\num{0.34}--\num{0.38} that disappears when the site is no longer shared between
training and test.

The two-tool design of Section~\ref{sec:xai} is a limitation as much as a
feature. Because SHAP was read only on frequency and Grad-CAM only on time, the
study never obtains a single explanation spanning both axes, and the one
genuinely cross-method check available --- agreement among the four independent
SHAP frequency profiles --- came out weak (\num{0.514}). A time--frequency
attribution that could be trusted on both axes at once, for instance a
higher-resolution CAM or an input-space gradient method on the spectrogram, would
allow the frequency and timing claims to be tested jointly rather than
separately; we leave that to future work.

Several further limitations should be stated plainly. The evaluation rests on one
split of one corpus; recording-level splitting and joint stratification guard
against the most common leakage, but a single split cannot separate
representation effects from split-specific variance, and repeated splits or
nested cross-validation would give firmer footing. Both random forests overfit
their training folds and neither classical model is well calibrated. The
MFCC-to-frequency projection used for SHAP is sign-agnostic and approximate, and
disagrees with the exact PWP mapping. The CNN's Grad-CAM and Grad-CAM++ maps
agree at only \num{0.392}, so the temporal result is specific to one variant. The
segmenter labels only \SI{66.5}{\percent} of test windows with sufficient
confidence, and nothing is known about whether the excluded third differs
systematically --- if the low-confidence windows are the noisy ones, they may
also be the ones where site cues matter most. Finally, the two smallest
sub-databases contribute 8 and 17 test recordings and carry essentially no
statistical weight.

We also note what did work. All five models substantially exceed trivial
baselines, and the ranking of feature domains is stable and interpretable: MFCC
outperforms PWP under both classifier families, and the log-Mel CNN attains the
best ranking quality (AUC \num{0.961}) despite an operating point that gives it
no MAcc advantage. The comparison the study was designed to make is answered,
even though the answer is not the one anticipated.
